% ===================================================================
% CHAPTER 7: SECOND-ORDER PRIMITIVES -- COHERENT STRUCTURES
% ===================================================================
\chapter{Second-Order Primitives: Coherent Structures}
\label{ch:second_order}

% Ported from: Separatrix_and_OW_Paper/Draft_11.tex (the Separatrix/
% Okubo-Weiss paper), which supersedes the older port of this material
% in thesis.tex Chapter 5. Draft_11 is treated as the primary source
% throughout this chapter; thesis.tex is used only in the few places
% noted below where it carries description Draft_11 dropped.

\section{Introduction and Related Work}
% New text (opening framing) followed by material ported from
% Draft_11.tex Sec. I (Introduction), trimmed to what is specific to
% coherent structures. Shared adaptive-navigation and scalar-field
% literature is in Chapter 1.

This chapter is the second rung of the fit-order ladder introduced in
Chapter~\ref{ch:intro}. The first rung (Chapter~\ref{ch:first_order}) fits
a plane to each velocity component from three robots, which gives the
velocity gradient at one point. A linear fit gives the velocity gradient
at one point; a second-order fit gives it as a field over the formation,
so any quantity built from it comes with its own gradient in closed form.
A single-point gradient can locate and classify an isolated critical
point, but it carries no information about how the field varies away from
that point, so it cannot find or follow a curve. This chapter fits a full
quadratic to each velocity component from six robots, recovering the
gradient and the Hessian from one synchronized sample, and uses the
result to build two adaptive navigation primitives that acquire and ride
a coherent structure from instantaneous measurements alone.

Formation-based Hessian recovery is established for scalar fields: a ring
of five robots plus one at the center recovers a scalar field's gradient
and Hessian in closed form from one synchronized sample \cite{15}. This
chapter extends that formation to a two-component field, recovering the
Jacobian and its second derivatives together, which is what makes the
surrogate fields of Section~\ref{sec:ch7:surrogates} available in closed
form.

Beyond critical points, vector fields have curves that are useful for
navigation. Dynamic ocean flows contain naturally occurring boundaries
known as coherent structures that influence where floating materials like
debris and larvae will accumulate or separate \cite{10,11}.
Mapping these exact curves allows teams to strategically deploy resources
in highly targeted search areas \cite{11,12,13,14}.
Without adaptive navigation, a large scale team must sample a much larger
basin to determine where these structures lie.

Coherent structures are typically identified through two primary methods.
The first relies on Lagrangian coherent structures (LCS) \cite{47},
typically computed as ridges in the finite-time Lyapunov exponent field
\cite{49}, with ongoing research addressing the Lagrangian/Eulerian
tradeoff and sparse-trajectory gradient recovery
\cite{48,52,50,51}. The second method uses instantaneous
Eulerian criteria, such as the Okubo-Weiss parameter, to partition flows
based solely on the local velocity gradient \cite{43,44}. This
instantaneous approach has been generalized to three dimensions
\cite{45}, refined to second order \cite{46}, and applied at scale
in mesoscale eddy censuses \cite{7}. However, both of these methods
characterize these structures offline using a reconstructed field, which
is not well suited to real-time adaptive navigation of a field that
evolves in time.

New work is being done on using adaptive navigation techniques with
multirobot systems in these scenarios. Michini et al.\ bracket a presumed
stable or unstable manifold with three robots using only local velocity
measurements \cite{34}, an approach inspired by the PIM triple
procedure, extended to $N$ robots \cite{35} and validated
experimentally in gyre-like flows \cite{36}. The robots begin on the
manifold and are advected by the flow, and the manifold's identity is
supplied in advance. The saddle point is treated by neither the
controller nor its analysis, and that work reports veering away from the
structure near a saddle in flow-tank experiments. Kularatne and Hsieh
confirm the boundary by computing the local FTLE field on board, at the
cost of a nonzero integration horizon and a commitment to attracting
structures \cite{1}. In both cases, the robotic formation
pre-straddles the structure, and does not acquire a structure it does not
already occupy.

This chapter shows that a second-order fit of the local flow is enough to
acquire and track such a structure from instantaneous measurements. The
contributions are
\begin{itemize}
\item a cooperative estimator that recovers the local quadratic model of a
two-component field from one synchronized set of measurements across six
robots, the minimum for the unconstrained model, and that fails only when
the robots lie on a common conic;
\item two surrogates read in closed form from that fit, the
velocity-gradient determinant $D$ and the smaller strain eigenvalue $s_1$
\cite{52};
\item two adaptive navigation primitives, the $D$ tracker and the $s_1$
tracker, that navigate to the structure and ride along it without
pre-straddling;
\item a tradeoff between them, in which the $D$ tracker tolerates more
measurement noise while the $s_1$ tracker, being objective, holds its path
under a rotating observer, which matters for a cluster whose heading is
dead-reckoned rather than measured.
\end{itemize}
Both primitives are demonstrated on a simulated double gyre and on
measured Santa Barbara Channel radar currents.

\section{Six-Robot Quadratic Estimation}
\label{sec:ch7:estimation}
% Ported from: Draft_11.tex Sec. II-A (Local Vector Field Polynomials)
% and Sec. II-D (Determinant Field, Strain Eigenvalue Field), specialized
% to the six-robot quadratic case. The general monomial-basis estimator,
% minimality proof, and common-conic degeneracy result are in Chapter 4
% (Sec. 4.3); only the second-order-specific closed forms are repeated
%% NOTE(new-math, checked 2026-09-29): eq:ch7:grad_D is not written in Draft_11.
%% Checked against eq:ch7:Jhat (u_xx=a5, u_xy=a4, u_yy=a6): both rows match.
% here. The gradient of $D$ (eq. \ref{eq:ch7:grad_D}) is a new derivative
% of Draft_11's eq. \ref{eq:ch7:Jhat}, written out because Draft_11 never
% states it explicitly (it is used only implicitly through the
% orthogonal-gradient argument of Sec.~\ref{sec:ch7:d_definition}); the
% same closed form appears, in $u_x, u_y, \ldots$ notation, in
% thesis.tex's "Closed-Form Recovery of $D$, $\nabla D$, $\mathbf{H}_D$."

Chapter~\ref{ch:estimation} states the general local-polynomial
estimator, proves that six robots are necessary and (generically)
sufficient for the unconstrained planar quadratic model, and
characterizes the resulting degeneracy as six points lying on a common
conic. This section specializes that estimator to give the closed-form
quantities the two trackers of this chapter consume.

A planar vector field assigns a velocity $\mathbf{v}(\mathbf{p}) =
(u(\mathbf{p}), v(\mathbf{p}))$ to each position $\mathbf{p} = (x, y)$.
Coordinates are measured from the formation centroid $\mathbf{p}_c$
throughout this chapter, so $\mathbf{p}_c = \mathbf{0}$. The second-order
Taylor expansion about the centroid is
\begin{align}
    u(\mathbf{p}) &= a_1 + a_2 x + a_3 y
        \nonumber \\
        &\quad + a_4 xy + \tfrac{1}{2} a_5 x^2 + \tfrac{1}{2} a_6 y^2
        + \mathcal{O}(\|\mathbf{p}\|^3),
        \label{eq:ch7:quad_model_u} \\
    v(\mathbf{p}) &= b_1 + b_2 x + b_3 y
        \nonumber \\
        &\quad + b_4 xy + \tfrac{1}{2} b_5 x^2 + \tfrac{1}{2} b_6 y^2
        + \mathcal{O}(\|\mathbf{p}\|^3),
        \label{eq:ch7:quad_model_v}
\end{align}
where $\mathcal{O}(\|\mathbf{p}\|^3)$ is the truncation error. The twelve
coefficients are constants, each a derivative of the field at the
centroid: $a_1 = u(\mathbf{p}_c)$, $a_2 = \partial u/\partial
x|_{\mathbf{p}_c}$, $a_3 = \partial u/\partial y|_{\mathbf{p}_c}$, $a_4 =
\partial^2 u/\partial x \partial y|_{\mathbf{p}_c}$, $a_5 = \partial^2
u/\partial x^2|_{\mathbf{p}_c}$, $a_6 = \partial^2 u/\partial
y^2|_{\mathbf{p}_c}$, and likewise for $b_1, \ldots, b_6$ with $v$ in
place of $u$. The one-half factors cancel the factor of 2 that
differentiation introduces. The zeroth-order coefficients give the fitted
flow at the centroid, $\hat{\mathbf{v}}_0 = (\hat{a}_1, \hat{b}_1)$, used
by both trackers below as the measured local velocity.

With robot positions $\mathbf{p}_i = (x_i, y_i)$, define the quadratic
basis vector
\begin{equation}
    \boldsymbol{\phi}(\mathbf{p}) =
    \bigl[\,1, \; x, \; y, \; xy, \;
    \tfrac{1}{2}x^2, \; \tfrac{1}{2}y^2\,\bigr]^\top,
\end{equation}
whose entries are the terms of
(\ref{eq:ch7:quad_model_u})--(\ref{eq:ch7:quad_model_v}) in order.
Stacking the six basis vectors row-wise gives the $6 \times 6$ formation
matrix $\boldsymbol{\Phi}$, and the coefficient estimates
$\hat{\boldsymbol{\theta}}_u = [\hat{a}_1, \ldots, \hat{a}_6]^\top$ and
$\hat{\boldsymbol{\theta}}_v = [\hat{b}_1, \ldots, \hat{b}_6]^\top$ solve
\begin{equation}
    \boldsymbol{\Phi}\,\hat{\boldsymbol{\theta}}_u = \mathbf{m}_u, \qquad
    \boldsymbol{\Phi}\,\hat{\boldsymbol{\theta}}_v = \mathbf{m}_v,
    \label{eq:ch7:solve}
\end{equation}
where $\mathbf{m}_u = [u_1, \ldots, u_6]^\top$ and $\mathbf{m}_v = [v_1,
\ldots, v_6]^\top$ collect the velocity components sampled by robot $i$.
Six robots suffice provided $\boldsymbol{\Phi}$ is nonsingular, which
fails exactly when the robots lie on a common conic (Chapter
\ref{ch:estimation}); a standard six-robot ring is degenerate at every
radius, since the ring itself is the conic.

For this study we chose the pentagon-plus-center formation and use it
throughout. It is the smallest uniform ring plus center on which
\cite{15} recovers a scalar Hessian exactly for quadratic signals,
independent of the formation's orientation. The same symmetry makes the
noise gains of the fit here independent of formation heading
(Chapter~\ref{ch:estimation}).

\paragraph{Closed-form recovery of $D$, $\nabla D$, and $\mathbf{H}_D$.}
The fitted coefficients define the velocity gradient at every point of
the footprint. Differentiating the quadratic model gives a Jacobian whose
entries are affine in position,
\begin{equation}
    \hat{\mathbf{J}}(\mathbf{p}) =
    \begin{bmatrix}
        \hat{a}_2 + \hat{a}_5 x + \hat{a}_4 y &
        \hat{a}_3 + \hat{a}_4 x + \hat{a}_6 y \\
        \hat{b}_2 + \hat{b}_5 x + \hat{b}_4 y &
        \hat{b}_3 + \hat{b}_4 x + \hat{b}_6 y
    \end{bmatrix}.
    \label{eq:ch7:Jhat}
\end{equation}
The determinant $\hat{D}(\mathbf{p}) = \det\hat{\mathbf{J}}(\mathbf{p})$
is therefore exactly quadratic over the footprint, and its value,
gradient, and Hessian at the centroid are polynomials in the
coefficients. The value is $\hat{D}_0 = \hat{a}_2\hat{b}_3 -
\hat{a}_3\hat{b}_2$. Differentiating (\ref{eq:ch7:Jhat}) once more gives
the gradient,
\begin{equation}
    \hat{\mathbf{g}}_{D,0} = \nabla\hat{D}|_0 =
    \begin{bmatrix}
        \hat{a}_5\hat{b}_3 + \hat{a}_2\hat{b}_4
            - \hat{a}_4\hat{b}_2 - \hat{a}_3\hat{b}_5 \\
        \hat{a}_4\hat{b}_3 + \hat{a}_2\hat{b}_6
            - \hat{a}_6\hat{b}_2 - \hat{a}_3\hat{b}_4
    \end{bmatrix},
    \label{eq:ch7:grad_D}
\end{equation}
and the Hessian is
\begin{equation}
    \hat{\mathbf{H}}_{D,0} =
    \begin{bmatrix}
        2(\hat{a}_5\hat{b}_4 - \hat{a}_4\hat{b}_5) &
        \hat{a}_5\hat{b}_6 - \hat{a}_6\hat{b}_5 \\
        \hat{a}_5\hat{b}_6 - \hat{a}_6\hat{b}_5 &
        2(\hat{a}_4\hat{b}_6 - \hat{a}_6\hat{b}_4)
    \end{bmatrix},
    \label{eq:ch7:hess_det}
\end{equation}
which is constant over the footprint, since it is the Hessian of the
fitted $\hat{D}$ rather than of the true $D$. The Hessian of the true $D$
also carries products of $\mathbf{J}$ with third derivatives of the
field, which a quadratic fit cannot recover at any radius.

\paragraph{Closed-form recovery of the strain quantities.} The second
surrogate comes from the rate-of-strain tensor $\hat{\mathbf{S}} =
\tfrac{1}{2}(\hat{\mathbf{J}} + \hat{\mathbf{J}}^\top)$, the symmetric
part of the same fitted Jacobian. At the centroid its mean, deviatoric
normal, and shear strain are $\hat{\mu} = (\hat{a}_2 + \hat{b}_3)/2$,
$\hat{s}_n = (\hat{a}_2 - \hat{b}_3)/2$, and $\hat{s}_s = (\hat{a}_3 +
\hat{b}_2)/2$, and each is affine in position. The eigenvalues of
$\hat{\mathbf{S}}$ are $\hat{s}_{1,2} = \hat{\mu} \mp \hat{r}$ with
$\hat{r} = \sqrt{\hat{s}_n^2 + \hat{s}_s^2}$, and their eigenvectors
$\mathbf{e}_1$ and $\mathbf{e}_2$ are the compression and stretching
directions. The gradient of the smaller eigenvalue follows by the chain
rule,
\begin{equation}
    \nabla\hat{s}_1 = \nabla\hat{\mu}
        - \frac{\hat{s}_n\nabla\hat{s}_n + \hat{s}_s\nabla\hat{s}_s}{\hat{r}},
    \label{eq:ch7:grad_s1}
\end{equation}
with $\nabla\hat{\mu}$, $\nabla\hat{s}_n$, and $\nabla\hat{s}_s$ read from
the second-order coefficients. Unlike $\hat{D}$, $\hat{s}_1$ has no usable
Hessian: it is an affine function minus the norm of an affine map, so it
is concave over the footprint for every field and formation.

\section{Eulerian Surrogates}
\label{sec:ch7:surrogates}
% Ported from: Draft_11.tex Sec. II-D (Eulerian Surrogates), the
% argument for why D and s1 mark the separatrix. The closed-form
% equations these subsections reference are given in
% Sec.~\ref{sec:ch7:estimation} above rather than repeated here.

Lagrangian methods locate coherent structures by integrating particle
trajectories over a finite time window \cite{47,49}. This chapter
instead reads them from the instantaneous geometry of the fitted velocity
field, as Eulerian coherent structures \cite{52}. Two fields computed
from the six-robot fit mark the separatrix: the determinant of the
velocity gradient, and the smaller eigenvalue of its symmetric part.

\subsection{Determinant Field}

For incompressible flow the determinant $\hat{D}$ of
(\ref{eq:ch7:Jhat})--(\ref{eq:ch7:hess_det}) is a negative multiple of the
Okubo-Weiss parameter \cite{43,44}, so the two carry the same
partition of the flow. Where $D < 0$ the flow is strain-dominated, and
where $D > 0$ it is rotation-dominated. In the flows studied here,
separatrices run along the trenches of $D$ inside the strain-dominated
regions.

\subsection{Strain Eigenvalue Field}

Serra and Haller define objective Eulerian coherent structures (OECS)
from the eigenvalues and eigenvectors of $\hat{\mathbf{S}}$
\cite{52}. An attracting OECS is a trench of $s_1$ tangent to
$\mathbf{e}_2$, and floating material has been observed at sea to collect
onto such curves \cite{11}. A repelling OECS is a ridge of $s_2$
tangent to $\mathbf{e}_1$, along which neighboring material separates. In
incompressible flow $s_2 = -s_1$, so repelling structures are also
trenches of $s_1$, and one tracker rides both.

\subsection{Relation Between the Surrogates}

Both fields are read from the same fitted Jacobian, so moving between
them changes nothing in the estimator. Where the vorticity $\omega =
\partial v/\partial x - \partial u/\partial y$ vanishes, $\mathbf{J}$
equals $\mathbf{S}$ and $D = s_1 s_2$, which is $-s_1^2$ in incompressible
flow. Along such a curve the two fields share their trenches. Elsewhere
the spin $\mathbf{W} = \mathbf{J} - \mathbf{S}$, the antisymmetric part of
$\mathbf{J}$ that carries the vorticity, separates them. It also
separates them across observers. A frame rotating clockwise at rate
$\Omega$, with rotation $\mathbf{Q}$, adds a uniform spin that raises
$\omega$ to $\omega + 2\Omega$. The strain tensor only co-rotates, so
$s_1$ is objective \cite{52}, while the determinant keeps the spin,
\begin{equation}
    D' = D + \Omega\,\omega + \Omega^2 .
    \label{eq:ch7:D_not_objective}
\end{equation}
Relatively rotating observers therefore disagree about where the trenches
of $D$ lie. The double-gyre example of Chapter~\ref{ch:preliminaries}
gives both fields in closed form and shows where their trenches coincide
and where they part.

\section{$D$ Tracker}
\label{sec:ch7:d_tracker}

\subsection{Definition}
\label{sec:ch7:d_definition}
% Ported from: Draft_11.tex Sec. III-C (The D Tracker). The Table I
% operating points (Draft_11 Sec. III-A) are reproduced here since both
% trackers draw their parameters from it; the surrounding three-layer
% control architecture description is in Chapter 4.

Both trackers of this chapter share the operating point of
Table~\ref{tab:ch7:params}, fixed across all experiments in this chapter
unless noted.

\begin{table}[!t]
\caption{Operating points. The double-gyre column is non-dimensional,
with one time unit equal to one second of the Decabot's identified
dynamics. Section~\ref{sec:ch7:sb_setup} converts the ocean column to
physical units.}
\label{tab:ch7:params}
\centering
\begin{tabular}{llll}
\hline
Parameter & Double gyre & Ocean & Units \\
\hline
$c_{\max}$              & 0.04  & 0.033  & length/time \\
$k$                     & 3.0   & 1.94   & none \\
$\rho$                  & 0.075 & 0.098  & length (5.0 km) \\
$v_{\text{max,robot}}$  & 0.3   & 0.3    & length/time \\
$v_{\text{stiction}}$   & 0.025 & 0.002  & length/time \\
$\alpha_{\text{mom}}$   & 0.7   & 0      & none \\
$\varepsilon_{\text{raw}}$ & $10^{-3}$ & $10^{-3}$ & $1/\text{time}^2$ \\
$\varepsilon_{\dim}$    & 0.025 & 0.025  & length$^2$ \\
$g_\perp$               & 1.0   & 1.0    & length$^2$ \\
$s_{\mathrm{trim}}$     & 0.05  & 0.3    & 1/time \\
$g_{\text{capture}}$    & 0.15  & 0.15   & 1/(time$\cdot$length) \\
$r_{\text{band}}$       & 0.05  & 0.05   & 1/time \\
\hline
\end{tabular}
\end{table}

On the double-gyre benchmark of Chapter~\ref{ch:preliminaries}, the $D$
landscape around the separatrix is a mountain pass. Two minima of $D$ sit
at the flow saddles, and the separatrix is the path of steepest descent
to each from a saddle point of $D$ at the origin, the crest. Around the
crest the landscape is a hyperbolic paraboloid, curving up across the
separatrix and down along it. The $D$ tracker drives the centroid onto
this trench and along it, over the crest, and down toward the next
minimum.

A trench is taken here in the height-ridge sense \cite{57}, a curve on
which the gradient is orthogonal to a Hessian eigenvector with positive
eigenvalue. The tracker reads this frame from $\hat{\mathbf{H}}_{D,0}
\mathbf{w}_j = \lambda_j \mathbf{w}_j$, with $\|\mathbf{w}_j\| = 1$,
$\lambda_1 \leq \lambda_2$, $\mathbf{w}_2$ across the trench, and
$\mathbf{w}_1$ along it, signed by $\hat{\mathbf{v}}_0^\top \mathbf{w}_1
\geq 0$ to point downstream. The identification is exact on the
double-gyre benchmark, where $\mathbf{H}_D$ is diagonal, and approximate
in general. The primitive assumes the true surface has the shape of the
crest, convex across and concave along; where that assumption fails
(Chapter~\ref{ch:preliminaries}) the fit does not register the change.

A full Newton step $-\hat{\mathbf{H}}_{D,0}^{-1}\hat{\mathbf{g}}_{D,0}$
leads to the stationary point of the fitted hyperbolic paraboloid, which
would stop the cluster at the crest. The tracker keeps only its
across-trench component and sets the along-trench speed separately, from
the along-trench gradient away from the crest and from the measured flow
near it, where that gradient vanishes. The flow takes over inside a
shallow-determinant band $\mathcal{B}$, the set of points where
\begin{equation}
    |\hat{D}_0| < \varepsilon_{\mathrm{raw}}
    \quad\text{or}\quad
    \frac{|\hat{D}_0|}{\|\hat{\mathbf{H}}_{D,0}\|_F}
        < \varepsilon_{\dim}.
    \label{eq:ch7:band}
\end{equation}
The second test measures depth against curvature, in units of length
squared, so the band scales with the landscape. The band surrounds the
zero set of $D$, which the separatrix meets at the crest and the cluster
may also cross during acquisition, where a trigger changes only the
along-trench speed. The along-trench speed is
\begin{equation}
    v_\parallel =
    \begin{cases}
        \hat{\mathbf{v}}_0^\top \mathbf{w}_1,
            & \mathbf{p}_c \in \mathcal{B}, \\[6pt]
        \dfrac{|\mathbf{w}_1^\top \hat{\mathbf{g}}_{D,0}|}{|\lambda_1|},
            & \text{otherwise},
    \end{cases}
    \label{eq:ch7:vpar}
\end{equation}
and the command, scaled by a navigation gain $k$, with each channel
saturated by $\mathrm{sat}(c) = c_{\max}\tanh(c/c_{\max})$, is
\begin{equation}
    \dot{\mathbf{p}}_c =
    \underbrace{\mathrm{sat}(v_\parallel)\,\mathbf{w}_1}_{\text{along the trench}}
    + \underbrace{\mathrm{sat}\!\left(-\frac{\mathbf{w}_2^\top
      \hat{\mathbf{g}}_{D,0}}{\lambda_2}\right)\mathbf{w}_2}_{\text{across the trench}}.
    \label{eq:ch7:d_law}
\end{equation}
The saturation caps the centroid speed at $\sqrt{2}\,k\,c_{\max}$. Both
branches of (\ref{eq:ch7:vpar}) are nonnegative, so the along-trench
motion never opposes the flow.

When $\hat{\mathbf{H}}_{D,0}$ is definite the trench frame is undefined.
The tracker then reverts to the signed Newton step
$\sum_j \mathrm{sat}(-\mathbf{w}_j^\top\hat{\mathbf{g}}_{D,0}/
\lambda_j)\mathbf{w}_j$ off the band and to the saturated flow on it. On a
positive-definite fit that step settles at the fitted minimum, so the
same sign test is the capture condition,
\begin{equation}
    \lambda_1 \lambda_2 \geq 0,
    \label{eq:ch7:d_capture}
\end{equation}
which the true $D$ meets at a flow saddle, where it has a local minimum.
A negative-definite fit carries no stability claim. On the benchmark the
test fires only near the flow saddles and domain corners, where second
derivatives vanish and truncation error sets the sign of the fitted
eigenvalues.

\subsection{Transverse Stability}
\label{sec:ch7:d_stability}
% Ported from: Draft_11.tex Sec. III-C (closing paragraph) and the
% Appendix ("Stability of the Trackers"), statement only. The proof is
% ported to Appendix~\ref{app:stability} by another pass.

The transverse channel of (\ref{eq:ch7:d_law}) reduces to $\dot{n} =
-k\,\mathrm{sat}(a_\perp n)$, where $n$ is the distance across the trench
and $a_\perp = \kappa_\perp/\hat{\kappa}_\perp$ is the ratio of true to
fitted transverse curvature. For $a_\perp > 0$ the cluster converges to
the trench, exponentially once $|n| < c_{\max}/a_\perp$, and estimation
error bounded by $\delta$ leaves it practically stable at $\limsup |n|
\leq \delta/a_\perp$. Since $v_\parallel \geq 0$, traversal never
reverses.

\begin{proposition}
\label{prop:ch7:transverse}
Let $\dot{n} = -k\,\mathrm{sat}(a_\perp n + e)$ with $a_\perp > 0$ and
$|e| \leq \delta$. Then $\limsup_{t\to\infty}|n| \leq \delta/a_\perp$. If
$e = 0$, $|n|$ falls below $c_{\max}/a_\perp$ in finite time, then decays
at rate at least $k\,a_\perp\tanh 1$.
\end{proposition}

For the $D$ tracker, writing the fitted transverse gradient as
$\hat\kappa_\perp n$ and the true one as $\kappa_\perp n + \epsilon$ with
$|\epsilon| \leq \delta_g$ gives $a_\perp = \kappa_\perp/\hat\kappa_\perp$
and $e = \epsilon/\hat\kappa_\perp$ in the proposition above, so the $D$
tracker additionally requires $\hat\kappa_\perp > 0$, that is, the fitted
Hessian must agree in sign with the true one across the trench. The full
derivation and proof are in Appendix~\ref{app:stability}.

\subsection{Simulation Verification}
\label{sec:ch7:d_verification}
% Ported from: Draft_11.tex Sec. IV-B (Clean Runs), the D-tracker-
% specific terminal-step result. The shared setup and acquisition/
% traversal narrative is in Section~\ref{sec:ch7:dg_bench}.

\begin{figure}[h]
  \centering
  \includegraphics[width=0.75\textwidth]{ch7_traverse_vs_logic_c.png}
  % FIGURE SOURCE: Paper_Writing/Separatrix_and_OW_Paper/figures/traverse_vs_logic_c.png
  \caption{$s_1$ tracker versus $D$ tracker from six matched starts
  ($\bullet$), noise-free. Both traverse the same network through the
  isotropic point at the origin. The $s_1$ tracker captures at the far
  saddle ($\times$), while the $D$ tracker continues onto the crossing
  wall trench past each saddle it reaches, ending at the square markers.}
  \label{fig:ch7:d_tracker_sim}
\end{figure}

The $D$ tracker passes within 0.02 of $\mathbf{p}^*_1 = (0, -0.5)$ in all
six runs of Fig.~\ref{fig:ch7:d_tracker_sim} and continues onto the
crossing wall trench. Near the saddle the fitted curvature is set by
truncation, so (\ref{eq:ch7:d_capture}) fires only intermittently
(Section~\ref{sec:ch7:d_definition}).

\section{$s_1$ Tracker}
\label{sec:ch7:s1_tracker}

\subsection{Definition}
\label{sec:ch7:s1_definition}
% Ported from: Draft_11.tex Sec. III-D (The s_1 Tracker).

The $s_1$ landscape has the same two minima as the $D$ landscape, and
between them the separatrix is again a trench. Where the $D$ trench
crosses a mountain pass, the $s_1$ trench crosses a sharp fold. In
incompressible flow $s_1 = -r \leq 0$, so the trench rises at most to
$s_1 = 0$, where the strain vanishes. Since $D = -s_1^2$ along the
separatrix (Section~\ref{sec:ch7:surrogates}), squaring rounds the fold
into the smooth crest of $D$.

The fitted $\hat{s}_1$ is concave for every field and formation
(Section~\ref{sec:ch7:estimation}), so unlike $\hat{D}$ it never shows the
upward curvature a trench has across it, and no Newton step is available.
Both channels are built instead from $\nabla\hat{s}_1$ and the strain
eigenframe, with $\hat{r}$ measuring how well resolved that eigenframe
is.

The trench tangent is a strain eigenvector, $\mathbf{e}_2$ on an
attracting segment and $\mathbf{e}_1$ on a repelling one, so no fixed
choice works across the isotropic point of $\mathbf{S}$. On the trench
the transverse derivative of $s_1$ vanishes, so away from a stationary
point $\nabla\hat{s}_1$ lies almost along the trench, and the tangent is
the eigenvector it is most nearly parallel to, signed against a
reference,
\begin{equation}
    \mathbf{t}_{\mathrm{raw}} =
    \arg\max_{\mathbf{e} \in \{\mathbf{e}_1, \mathbf{e}_2\}}
    \bigl\lvert \nabla\hat{s}_1^\top \mathbf{e} \bigr\rvert,
    \qquad
    \mathbf{t} = \mathrm{sgn}\bigl(\mathbf{t}_{\mathrm{ref}}^\top
    \mathbf{t}_{\mathrm{raw}}\bigr)\,\mathbf{t}_{\mathrm{raw}},
    \label{eq:ch7:tangent_select}
\end{equation}
where $\mathbf{t}_{\mathrm{ref}}$ is the previous step's tangent, or
$\hat{\mathbf{v}}_0$ on the first ride step, since the measured flow is
not objective (Section~\ref{sec:ch7:rotating}). The rule swaps from
$\mathbf{e}_2$ to $\mathbf{e}_1$ at the isotropic point with no dedicated
test.

\begin{figure}[h]
  \centering
  \includegraphics[width=\textwidth]{ch7_s1_channels.png}
  % FIGURE SOURCE: Paper_Writing/Separatrix_and_OW_Paper/figures/s1_channels.png
  \caption{The two channels of the $s_1$ tracker, on the benchmark trench.
  (a) Tangent selection, (\ref{eq:ch7:tangent_select}). At each sample the
  strain eigenframe is drawn, the selected eigenvector heavy and the
  rejected one light, with $\nabla\hat{s}_1$ in orange. The argmax takes
  $\mathbf{e}_2$ on the attracting segment and $\mathbf{e}_1$ on the
  repelling segment, swapping through the isotropic point. (b) The
  command. Travel along $\mathbf{t}$ is open loop at cruise speed, while
  $\mathbf{P}$ removes the along-trench component of the gradient so
  descent acts only across the trench.}
  \label{fig:ch7:s1_channels}
\end{figure}

Until $\hat{s}_1$ first falls below $-s_{\mathrm{trim}}$ the command is
gradient descent, $-\mathrm{sat}(g_\perp\nabla\hat{s}_1)$, with transverse
gain $g_\perp$. From then on the tracker rides,
\begin{equation}
    \dot{\mathbf{p}}_c =
    c_{\max}\tanh(1)\,\mathbf{t}
    - \mathrm{sat}\!\bigl(g_\perp\,\mathbf{P}\,\nabla\hat{s}_1\bigr),
    \qquad
    \mathbf{P} = \mathbf{I} - \mathbf{t}\mathbf{t}^\top,
    \label{eq:ch7:s1_law}
\end{equation}
scaled by $k$ as in (\ref{eq:ch7:d_law}), where $\mathrm{sat}$ saturates a
vector's magnitude and preserves its direction (Fig.
\ref{fig:ch7:s1_channels}). The projector removes the along-trench
component of the gradient, so the cluster descends only across the trench
while riding along it at speed $k\,c_{\max}\tanh 1$. The latch depth
$s_{\mathrm{trim}}$ is a small fraction of the benchmark's well depth
$-\pi^2 A \approx -0.987$ (Chapter~\ref{ch:preliminaries}), so from most
starts the ride begins at the first step and projected descent carries
the cluster onto the trench. Below a strain floor $r_{\mathrm{band}}$ the
eigenframe is unresolved. If $\hat{r} < r_{\mathrm{band}}$ before any
tangent exists, (\ref{eq:ch7:tangent_select}) has no direction, and the
command rides $\hat{\mathbf{v}}_0/\|\hat{\mathbf{v}}_0\|$ with $\mathbf{P}
= \mathbf{I}$. Only the first ride step can reach this case.

The tracker also stops at a well. Its capture test is built on the
gradient, since $\hat{\mathbf{H}}_{s_1}$ is negative semidefinite for
every field and formation (Section~\ref{sec:ch7:estimation}) and cannot
certify a minimum. The fitted gradient tracks the true gradient even
where the fitted curvature carries the wrong sign, and a flow saddle is a
true two-dimensional minimum of $s_1$, so a vanishing gradient suffices.
When $\hat{r} > r_{\mathrm{band}}$, $\|\nabla\hat{s}_1\| <
g_{\mathrm{capture}}$ for a capture threshold $g_{\mathrm{capture}}$
(Table~\ref{tab:ch7:params}), and $\hat{s}_1 < -4s_{\mathrm{trim}}$, the
tracker drops the ride term and descends $\hat{s}_1$, holding the cluster
at the minimum until $\hat{s}_1$ rises back above $-s_{\mathrm{trim}}$.

\subsection{Transverse Stability}
\label{sec:ch7:s1_stability}
% Ported from: Draft_11.tex Sec. III-D (closing paragraph) and the
% Appendix ("Stability of the Trackers"), statement only. The proof is
% ported to Appendix~\ref{app:stability} by another pass.

The transverse channel of (\ref{eq:ch7:s1_law}) reduces to the same form
as Proposition~\ref{prop:ch7:transverse}, with $a_\perp =
g_\perp\kappa_\perp$ set by the field rather than normalized by a fitted
curvature, and $e = g_\perp\epsilon$. Both bounds equal
$\delta_g/\kappa_\perp$, independent of gain and fitted curvature, but
only the $D$ tracker needs $\hat\kappa_\perp > 0$. Where $\kappa_\perp$
vanishes over a length, as at the fold, the ride carries the cluster
across (Appendix~\ref{app:stability}). A frame rotating at $\Omega$
perturbs $\dot n$ by at most $\Omega\lVert\mathbf{p}\rVert$, leaving $n$
bounded while $\Omega\lVert\mathbf{p}\rVert < k\,c_{\max}$
(Section~\ref{sec:ch7:rotating}). The ride direction is seeded once from
the measured flow and is guaranteed only under exact estimates.

\subsection{Simulation Verification}
\label{sec:ch7:s1_verification}
% Ported from: Draft_11.tex Sec. IV-B (Clean Runs), the s1-tracker-
% specific terminal-step result. The shared setup and acquisition/
% traversal narrative is in Section~\ref{sec:ch7:dg_bench}.

\begin{figure}[h]
  \centering
  \includegraphics[width=0.75\textwidth]{ch7_separatrix_formation.png}
  % FIGURE SOURCE: Paper_Writing/Separatrix_and_OW_Paper/figures/separatrix_formation.png
  \caption{The six robots from start S1 (star) under both trackers,
  zoomed to the separatrix $x = 0$ (dashed). Thin lines are robot paths
  and the heavy line is the centroid. Outlines mark the
  pentagon-plus-center formation at the start, during the ride, and at
  the far saddle ($\times$).}
  \label{fig:ch7:s1_tracker_sim}
\end{figure}

The $s_1$ tracker settles within 0.011 to 0.022 of $\mathbf{p}^*_1 = (0,
-0.5)$ in all six runs, capturing at the far saddle rather than
continuing onto the crossing wall trench, as shown for start S1 in
Fig.~\ref{fig:ch7:s1_tracker_sim}.

\section{Behavior on the Double-Gyre Benchmark}
\label{sec:ch7:dg_bench}
% Ported from: Draft_11.tex Sec. IV-A (Setup) and IV-B (Clean Runs).
% Renamed from the skeleton's "Estimator Accuracy on the Double Gyre";
% the content Draft_11 places here is acquisition and traversal
% behavior under clean measurements, not an open-loop estimator-error
% sweep (estimator_accuracy_vs_noise.png discarded by the author).

\subsection{Setup}

The double-gyre experiments use the steady field of
Chapter~\ref{ch:preliminaries} ($A = 0.1$) on $[-1, 1] \times [-0.5,
0.5]$. All trials use the nominal formation with ring radius $\rho =
0.075$; trials vary only the initial heading, drawn uniformly from $[0,
2\pi)$ in the Monte Carlo sweeps of Sections~\ref{sec:ch7:noise} and
\ref{sec:ch7:rotating}, and zero in the noise-free runs below.

\subsection{Clean Runs}

These runs test acquisition, traversal, and the terminal step for both
primitives. On the double gyre the two fields share every trench
(Chapter~\ref{ch:preliminaries}). Clean trials run 400 to 600 steps and
stop at domain exit or, for the $D$ tracker, 150 steps after first saddle
contact.

Both trackers ran noise-free from six matched starts
(Fig.~\ref{fig:ch7:d_tracker_sim}), spanning both gyre halves, the
origin, and points up to $2.7\rho$ off the structure, acquiring the
network from all six in 0 to 13 steps for the $D$ tracker and 0 to 14 for
the $s_1$ tracker. They then ride the same segments in the same direction
and cross the isotropic point at the origin, with the formation
straddling the separatrix throughout (Fig.~\ref{fig:ch7:s1_tracker_sim}).
A run acquiring above the origin therefore rides an attracting OECS to
that point and continues onto a repelling one, on the same primitive and
the same scalar trench throughout. Neither needs a dedicated rule at the
crossing, and the $s_1$ tracker rides through the origin on ordinary
tangent selection. The primitives part only at the terminal step
(Sections~\ref{sec:ch7:d_verification} and~\ref{sec:ch7:s1_verification}).

\section{Behavior Under Noise}
\label{sec:ch7:noise}
% Ported from: Draft_11.tex Sec. IV-C (Behavior Under Noise).

\begin{figure}[h]
  \centering
  \includegraphics[width=0.75\textwidth]{ch7_flip_resolution.png}
  % FIGURE SOURCE: Paper_Writing/Separatrix_and_OW_Paper/figures/flip_resolution.png
  \caption{Far-saddle success of both trackers from the straddling start
  $(0, 0.35)$, $10^4$ trials per point, with each tracker's straddle
  retention (counted on successful trials) dashed. The dotted line marks
  50\% and the ticks each tracker's crossing. (a) Against $\sigma_{uv}$,
  $\sigma_p = 0$. (b) Against $\sigma_p$, $\sigma_{uv} = 0$.}
  \label{fig:ch7:flip_resolution}
\end{figure}

Both primitives sweep measurement and position noise separately from one
straddling start, $(0, 0.35)$, at 10\,000 trials per noise level from
independent headings. Noise-sweep trials run up to 500 steps and stop at
task completion or domain exit. Success is contact within 0.06 of the far
saddle $\mathbf{p}^*_1 = (0, -0.5)$, which a noise-free run from this
start reaches under either primitive, without formation collapse (an RMS
error above $4\rho$ in the fifteen inter-robot distances); each trial also
records straddle retention. A trial straddles when the robots fall on
both sides of the separatrix, and retention requires it never stop once
started, the failure \cite{34} reports. Holding one target fixed makes
a settle at the near saddle a directional failure rather than a loss of
the structure. Both succeed in every noise-free trial. Noise favors the
$D$ tracker, whose 50\% crossing is about four times the $s_1$ tracker's.

The $D$ tracker crosses 50\% at $\sigma_{uv} \approx 0.0077$, where
$\nabla\hat{D}$ stops being informative along the separatrix
($\sigma_{uv} \approx 0.0047$ to $0.0085$), and the $s_1$ tracker at
$\sigma_{uv} \approx 0.0019$ (Fig.~\ref{fig:ch7:flip_resolution}(a)).
Under position noise the ordering holds, $\sigma_p \approx 0.0081$
against 0.0023 (Fig.~\ref{fig:ch7:flip_resolution}(b)). The $s_1$
tracker's failure is a sign inversion at tangent selection
(Section~\ref{sec:ch7:s1_definition}). Supplying a clean gradient and
eigenframe to (\ref{eq:ch7:tangent_select}) at $\sigma_{uv} = 0.002$
restores the correct-saddle rate to 100\%. This fragility is attributed
to how long a sign decision persists before a new measurement checks it,
rather than to per-cycle estimation accuracy. The $D$ tracker re-signs
$\mathbf{w}_1$ against $\hat{\mathbf{v}}_0$ every cycle, while the $s_1$
tracker carries the sign in state through $\mathbf{t}_{\mathrm{ref}}$.
That reliance follows from objectivity and traversal together.
Objectivity rules out the fitted curvature
(Section~\ref{sec:ch7:surrogates}), the measured flow
(Section~\ref{sec:ch7:rotating}), and a world axis as sign references, and
traversal rules out the along-trench component of $\nabla\hat{s}_1$,
which reverses at the along-trench maximum the ride must cross, leaving
the primitive's own prior output. Straddle retention stays within five
points of success for the $D$ tracker but falls further below it for the
$s_1$ tracker, 47.6\% against 71.5\% at $\sigma_{uv} = 0.0015$. Most
$s_1$ trials that lose the straddle do so within 0.1 of the origin, where
the $s_1$ trench loses its transverse curvature
(Section~\ref{sec:ch7:s1_definition}).

\section{Behavior Under a Rotating Observer}
\label{sec:ch7:rotating}
% Ported from: Draft_11.tex Sec. IV-D (Behavior Under a Rotating
% Observer). This is the objectivity result and is reused in
% Chapter 8.

\begin{figure}[h]
  \centering
  \includegraphics[width=\textwidth]{ch7_objectivity_traverser.png}
  % FIGURE SOURCE: Paper_Writing/Separatrix_and_OW_Paper/figures/objectivity_traverser.png
  \caption{Objectivity trial, both primitives started from $(0.05, 0.40)$
  and run in the inertial frame (solid) and in a frame rotating at
  $\Omega = 0.23$ (dashed, pulled back to inertial coordinates). (a) The
  $s_1$ tracker's two runs coincide and capture at the same saddle. (b)
  The $D$ tracker's rotating-frame run loses the straddle across the
  crest (inset), the artifact predicted by (\ref{eq:ch7:D_not_objective}),
  and rejoins near the far saddle.}
  \label{fig:ch7:oecs_objectivity}
\end{figure}

Both primitives were run from $(0.05, 0.40)$, just inside the upper
saddle, on the same physical double gyre in two observer frames, the
inertial one and a frame rotating clockwise at $\Omega = 0.23$, a heading
drift of about $8^\circ$/h at the ocean time unit of
Section~\ref{sec:ch7:sb_setup}, where the field is
$\mathbf{v}'(\mathbf{p}', t) = \mathbf{Q}\mathbf{v}(\mathbf{Q}^\top
\mathbf{p}') + \Omega\,(-y', x')^\top$. The rate keeps the structure's
apparent speed along the tracked path, $\Omega\lVert\mathbf{p}\rVert \leq
0.11$, inside the command authority $k\,c_{\max} = 0.12$, the condition
Appendix~\ref{app:stability} requires. Rotating-frame trajectories are
pulled back to inertial coordinates and compared with the inertial run of
the same primitive.

Between acquisition and the far saddle, the $s_1$ tracker stays within
0.016 of the separatrix in the rotating frame, against 0.010 in the
inertial frame, keeps the straddle throughout, and both runs capture at
the same bottom saddle (Fig.~\ref{fig:ch7:oecs_objectivity}). The $D$
tracker's rotating-frame run is displaced by up to 0.084, against 0.010,
and loses the straddle for 36 steps across the crest, with all six robots
on one side of the separatrix, the failure \cite{34} reports.
Sweeping the rate, the $D$ tracker first loses the straddle at $\Omega
\approx 0.21$, where $\Omega\lVert\mathbf{p}\rVert$ at the crest is a
sixth of the command authority.

The frame's solid-body vorticity shifts the $D'$ landscape
(\ref{eq:ch7:D_not_objective}) and adds a transport term to the measured
flow, so both terms of the $D$ tracker are steered by a frame artifact.
On the separatrix $\omega = 0$ but $\partial_x\omega$ is nonzero
(Chapter~\ref{ch:preliminaries}), so $\Omega\omega$ tilts $D'$ across the
trench, moving the trench most at the origin and least at the saddles, so
the run rejoins the separatrix near $\mathbf{p}_1^*$. The $s_1$ tracker's
tests are scalar inequalities in quantities a rotation leaves invariant,
so its command co-rotates with the frame, and its residual is the
frame-transport term bounded in Appendix~\ref{app:stability}. Its one
non-objective input is the tangent seed taken from the measured flow,
which fixes the ride direction, and the frames agree while
$\Omega\lVert\mathbf{p}_{\text{seed}}\rVert <
\lvert\hat{\mathbf{v}}_0^\top\mathbf{t}_{\text{raw}}\rvert$.

\section{Santa Barbara Channel Demonstration}
\label{sec:ch7:sb_channel}
% Ported from: Draft_11.tex Sec. V (Santa Barbara Channel Trial), full.
% One sentence of data-provenance detail (CORDC, THREDDS-Ocean, retrieval
% date) is added from thesis.tex's "Ocean HFR Simulation" (lines
% ~3958-3968), which describes it more fully than Draft_11 does; every
% numeric operating-point and result value below is Draft_11's, not
% thesis.tex's, since thesis.tex's ocean trial used a different
% formation scale and gain and is not comparable.

\subsection{Setup}
\label{sec:ch7:sb_setup}

The ocean experiment uses hourly surface-current frames for the Santa
Barbara Channel from the HFRNet U.S.\ West Coast Radial-to-Total Vector
product \cite{58}, on the 2~km grid, which is public. The product is
produced by the Scripps Coastal Ocean Observing Research and Development
Center (CORDC) HFRNet network and archived at the NOAA National Centers
for Environmental Information, retrieved from the NOAA THREDDS-Ocean
server. The record is the one \cite{34} uses, the same network over
the same 28-h window, 29 frames from May 16, 2012 08:00 GMT to May 17,
2012 12:00 GMT, so these results are interpretable against prior work on
it. Kularatne and Hsieh track a time-reversed attracting boundary in the
same channel's May 2012 currents \cite{1}. Frames are interpolated
linearly in time. World coordinates on $[-0.65, 0.65]^2$ map affinely onto
a region centered on $(34.2^{\circ}\mathrm{N}, 120.4^{\circ}\mathrm{W})$.

Physical numbers follow from one unit block, length $L_w = 51.4$~km on
both axes and time $6000$~s, with a control period $\Delta t = 0.1$ time
units, so one step advances the field by $600$~s and the $168$-step trial
covers the record. The commanded-speed cap $\sqrt{2}\,k\,c_{\max}$ is then
$0.78$~m/s, within reach of a small autonomous surface vessel but above a
buoyancy glider. With longitude scaled by $\cos(34.2^{\circ})$ the map is
a uniform dilation, so the zero sets, trenches, and eigenframe are the
physical ones. The FTLE reference uses the same constants, over a 24-h
field for Fig.~\ref{fig:ch7:sb_channel} and 12~h for a persistence check.

The pentagon is enlarged to $\rho = 0.098$, $1.3\times$ the nominal
radius and $5.0$~km at this site, so the footprint spans about five cells
of the 2~km grid rather than under four, and starts at heading
$335^\circ$. At one control cycle per ten minutes any vessel time
constant is negligible, so $\alpha_{\text{mom}} = 0$. The gain and radius
set when the cluster reaches the bifurcation of the evolving ridge, and
so which branch it meets. The latch depth $s_{\mathrm{trim}} = 0.3$ is
set from this field's own $s_1$ scale. No noise level is estimated, since
the exactly determined fit leaves no residual, so the noise model of
Section~\ref{sec:ch7:noise} is validated on the double gyre only.

\subsection{Results}
\label{sec:ch7:sb_results}

\begin{figure}[h]
  \centering
  \includegraphics[width=\textwidth]{ch7_ocean_progression_2km.png}
  % FIGURE SOURCE: Paper_Writing/Separatrix_and_OW_Paper/figures/ocean_progression_2km.png
  \caption{$D$ and $s_1$ tracker paths from a single shared start at (a)
  7, (b) 14, (c) 21, and (d) 28~h, each drawn up to that time, with dots
  marking the current positions. The background is the 24-h forward FTLE
  field computed offline from the same 2~km data and anchored at the
  record start. Both follow the dominant ridge from the channel entrance
  to the bifurcation, where the $D$ tracker turns onto the southeastern
  branch toward the middle island. The $D$ run stops at 24~h, when its
  first robot reaches the island.}
  \label{fig:ch7:sb_channel}
\end{figure}

The field is divergent and not vorticity-free, so $s_2 = -s_1$ does not
hold and the two surrogates need not share trenches. Both were run with
no changes to the primitives, at the ocean operating point of
Table~\ref{tab:ch7:params}, from a shared start north of the channel
entrance, $(34.38^{\circ}\mathrm{N}, 120.41^{\circ}\mathrm{W})$. Over the
28-h trial both descend through the channel entrance and track the same
ridge to the bifurcation. There the $D$ tracker turns onto the
southeastern branch, and its run stops after 24~h, when its first robot
reaches the coast of the middle island. The $s_1$ tracker is still
descending when the record ends.

Fig.~\ref{fig:ch7:sb_channel} places both trials against the 24-h
forward-time FTLE field, computed offline in the style of \cite{49}.
Both follow the dominant ridge from the channel entrance to the
bifurcation closely. Mean distance from the tracked path to the nearest
95th-percentile ridge point is 1.5~km for the $D$ tracker and 2.1~km for
the $s_1$ tracker. The grid is 2~km and the formation radius 5.0~km. That
ridge was available to neither, and FTLE recomputed at four anchor times
shows it persisting while the finer structure downstream evolves.

The trial exercises tangent selection where neither the eigenframe nor
the trench network is axis-aligned, which the double gyre cannot do. No
two-dimensional minimum of $s_1$ lies within reach, so the capture test
never engages.

\section{Summary}
\label{sec:ch7:summary}
% New text, summarizing this chapter's own results. Draft_11's
% Conclusion also carries the cross-cutting objectivity synthesis
% (Chapter 8's territory) and future-work items (Chapter 9's
% territory, per the settled decision that Draft_11 Sec. VI stays out
% of this chapter); neither is repeated here.

A six-robot pentagon-plus-center formation recovers the local quadratic
model of a two-component field from one synchronized sample, giving the
Jacobian, its determinant $D$, and the smaller rate-of-strain eigenvalue
$s_1$ in closed form, together with their gradients and, for $D$, its
Hessian (Section~\ref{sec:ch7:estimation}). Both surrogates mark the
separatrix as a trench (Section~\ref{sec:ch7:surrogates}), and the $D$
tracker and $s_1$ tracker built on them acquire and ride the structure
from instantaneous measurements alone, with no pre-straddling and no
integration horizon.

On the double-gyre benchmark both trackers acquire, traverse, and cross
the isotropic point identically under clean measurements, parting only at
the terminal saddle. Under noise the $D$ tracker withstands about four
times the measurement noise the $s_1$ tracker does, with a matching
margin under position noise. Under a rotating observer the ranking
reverses: above a heading drift of about $7^\circ$/h the $D$ tracker's
formation no longer straddles the structure, while the $s_1$ tracker,
being objective, keeps the straddle and returns the same material curve
in both frames. A cluster with an absolute heading reference should carry
the $D$ tracker; one whose heading is dead-reckoned should carry the
$s_1$ tracker. On Santa Barbara Channel currents both primitives followed
the same transport corridor to its bifurcation, at mean distances of 1.5
and 2.1~km from an FTLE ridge neither could see, on a field neither
trials nor the double-gyre benchmark can fully stand in for.

% If hardware verification for second-order primitives happens before
% defense (per the March-June 2027 window discussed with the advisor),
% add a Hardware Verification subsection to each tracker above and a
% corresponding \section{Hardware Verification} here before this
% Summary.
